% xpar 2.0 wire-format specification
% Copyright (C) 2022-2026 Kamila Szewczyk
\documentclass[10pt,a4paper]{report}
\usepackage[T1]{fontenc}
\usepackage[utf8]{inputenc}
\usepackage{amsmath,amssymb,array,booktabs,caption,longtable,microtype}
\usepackage[margin=25mm]{geometry}
\usepackage[hidelinks]{hyperref}
\usepackage{xcolor}
\usepackage{fancyvrb}
\usepackage{enumitem}
\setlist{nosep}
\setlength{\emergencystretch}{3em}
\providecommand{\tightlist}{\setlength{\itemsep}{0pt}\setlength{\parskip}{0pt}}
\newcounter{none}
\newcounter{specparagraph}[section]
\makeatletter
\@addtoreset{specparagraph}{chapter}
\@addtoreset{specparagraph}{subsection}
\@addtoreset{specparagraph}{subsubsection}
\makeatother
\newcommand{\specp}{\refstepcounter{specparagraph}\noindent
  \makebox[0pt][r]{\textsuperscript{\scriptsize[\thespecparagraph]}\hspace{0.7em}}}
\title{xpar Container Format 2.0}
\author{Kamila Szewczyk}
\date{August 2026}

\begin{document}
\maketitle
\tableofcontents

\chapter{Scope}\label{scope}

\specp This document specifies the xpar 2.0 file format. It describes in depth
the byte representation, coding geometry, integrity and authentication,
packet constraints, volume layouts, deterministic construction,
parser requirements, and conformance vectors.

\specp xpar 2.0 is by design incompatible with other competing standards, such
as PAR2, PAR3, and xpar 1.0. Notwithstanding, a conforming reader shall detect
an xpar 1.0 signature, shall report an unsupported version, and shall not
interpret the input as xpar 2.0.

\specp The xpar 1.0 signatures are the five-byte sequence \texttt{XP}, one
major-version byte in the range 0 through 1, one minor-version byte in the
range 0 through 9, and a final byte in
\texttt{1}, \texttt{2}, \texttt{3}, or \texttt{4}; the four-byte sequence
\texttt{XPAS}; and the four-byte sequence \texttt{XPAL}. Detection of these
signatures shall produce only an unsupported-version diagnostic. This document
does not specify decoding of any xpar 1.0 object.

\chapter{Normative references}\label{normative-references}

\specp The following documents are normatively referenced:
\begin{itemize}
\item BLAKE3 specification, version 1.3.3, 9 June 2025;
\item RFC 3720, Appendix B.4, CRC-32C, April 2004;
\item ISO/IEC 10646:2020, UTF-8 encoding form;
\item FIPS PUB 180-4, Secure Hash Standard, August 2015.
\end{itemize}

\specp A dated reference denotes only the stated edition. The parameter,
serialization, domain-separation, and truncation rules in this document shall
apply where a referenced specification permits alternatives.

\chapter{Terms, definitions, and conformance}\label{terms}

\section{Conformance}

\specp An implementation conforms to this specification only if every output
it identifies as xpar 2.0 satisfies all applicable ``shall'' and ``shall not''
requirements and it rejects every input for which this specification requires
rejection.

\specp In this specification, ``shall'' expresses a requirement, ``shall not''
expresses a prohibition, ``should'' expresses a recommendation, and ``may''
expresses permission. Text identified as a constraint is normative. Examples
do not override a requirement or constraint.

\section{Terms and definitions}

\begin{description}[style=nextline]
\item[entry] A manifest object denoting a regular file, directory, symbolic
link, or hard-link alias.
\item[extent] A half-open interval in the global set-stream address space.
\item[generation] One manifest state and the independently coded stream range
appended by that state.
\item[set] An ordered byte stream, a manifest mapping entries to stream
extents, integrity metadata, and zero or more recovery slices.
\item[slice] A coding unit of exactly $Z$ bytes. The final data slice includes
zero padding that is not part of the logical stream.
\item[volume] One file that contains xpar packets, bare data, or an armoured
archive as specified by its layout.
\end{description}

\chapter{Byte and arithmetic representation}\label{representation}

\specp Byte interval $[a,b)$ contains byte offsets $a$ through $b-1$.
$A\mathbin{\|}B$ denotes byte-string concatenation. $\mathrm{u}N\mathrm{le}(x)$
denotes the unsigned $N$-bit little-endian representation of $x$.
Arithmetic constraints use mathematical integers and shall be checked without
wraparound.

\specp Unless a field definition states otherwise, integers are unsigned and
little-endian, offsets are measured from the beginning of the containing body,
reserved fields shall be written as zero, and padding bytes shall be zero.

\end{document}
